This study demonstrated that the Elliptic filter topology provides the most effective front-end solution for isolating gravitational wave chirps in the presence of strong terrestrial interference. By modeling the observatory as a Linear Time-Invariant (LTI) system and applying frequency-domain deconvolution, the structure of a 30--300~Hz Linear Frequency Modulation (LFM) chirp was successfully recovered from empirical LIGO strain data. The results showed that the Elliptic filter achieved the highest performance, with a correlation of $\rho = 0.7275$ and a Peak Signal-to-Noise Ratio (PSNR) of 5.8647~dB, outperforming both Butterworth and Chebyshev models. This can be attributed to its steep transition band, which effectively suppressed the dominant 60~Hz mains interference. In contrast, the poor performance of the Chebyshev~II filter ($\rho = 0.0689$) highlights that non-linear phase behavior can significantly distort transient signals, emphasizing that phase preservation is as important as magnitude attenuation.

To further improve the accuracy of gravitational wave reconstruction, it is recommended to move from static filter designs toward more phase-consistent architectures, while also expanding the range of signal models considered. Although the Elliptic filter provides the required steepness, integrating Bessel filters should be explored due to their maximally flat group delay, which may help reduce the phase distortion observed in this study. In addition, implementing an Adaptive Notch Filter can allow the system to track and suppress time-varying power-line harmonics more effectively, while a cascaded $8^{th}$-order configuration may offer better stability in practical hardware implementations. Moreover, future work should not be limited to LFM chirps but should also include Hyperbolic and Exponential chirps. These signal types better represent the non-linear inspiral behavior of binary mergers and provide improved robustness to Doppler effects, making the system more adaptable to variations in chirp mass and orbital eccentricity. Lastly, incorporating wavelet-based denoising after the deconvolution process is recommended to reduce high-frequency artifacts introduced during spectral inversion, leading to a cleaner and more accurate signal reconstruction.